\documentclass[12pt]{article} %{{{
\usepackage{showlabels}
\usepackage{amsmath,amsthm,amsfonts,amssymb,color,enumerate,xfrac,tikz,comment, pgfplots,soul, mathtools}
\usepackage{mathrsfs}
% \usepateckage[colorlinks,citecolor=blue,urlcolor=red]{hyperref}



% -------------------------------------------------
% Environments
% -------------------------------------------------
\theoremstyle{plain}
\newtheorem{theorem}{Theorem}[section]
\newtheorem{corollary}[theorem]{Corollary}
\newtheorem{example}[theorem]{Example}
\newtheorem{proposition}[theorem]{Proposition}
\newtheorem{lemma}[theorem]{Lemma}


\theoremstyle{definition}
\newtheorem{remark}[theorem]{Remark}
\newtheorem{assumption}{Assumption}
\newtheorem{definition}[theorem]{Definition}


\numberwithin{equation}{section}

% -------------------------------------------------
% User defined commands
% -------------------------------------------------
\def\cB{\mathcal{B}}
\def\cD{\mathcal{D}}
\def\cE{\mathcal{E}}
\def\cF{\mathcal{F}}
\def\cH{\mathcal{H}}
\def\cL{\mathcal{L}}
\def\cN{\mathcal{N}}
\def\cP{\mathcal{P}}
\def\cS{\mathcal{S}}
\def\cV{\mathcal{V}}


\def\D{\mathbb{D}}
\def\bR{\mathbb{R}}
\def\bC{\mathbb{C}}
\def\e{\varepsilon}
\def\E{\mathbb E}
\def\H{\mathbb{H}}
\def\P{\mathbb P}
\def\Z{\mathbb Z}


\newcommand{\DRL}[2]{{}_t D_{#1}^{#2}}
\newcommand{\lMr}[3]{\:{}_{#1} #2_{#3}}
\newcommand{\Norm}[1]{\left|\left|  #1   \right|\right|}

\def\blue{\color{blue}}
\def\red{\color{red}}

\def\todo{{\blue(To continue from here.)}}
\topmargin -0.4in
\headsep 0.4in
\textheight 9.0in
\oddsidemargin 0.02in
\evensidemargin 0.15in
\textwidth 6.3in


\DeclareMathOperator*{\esssup}{ess\,sup}

\newcommand{\DCap}[1]{\partial^{#1}}
\newcommand{\R}{\mathbb{R}}
\newcommand{\N}{\mathbb{N}}
\newcommand{\W}{\dot{W}}
\newcommand{\ud}{\ensuremath{ \mathrm{d}} }
\newcommand{\Ceil}[1]{\left\lceil #1 \right\rceil}
\newcommand{\Floor}[1]{\left\lfloor #1 \right\rfloor}
\newcommand{\sgn}{\text{sgn}}
\newcommand{\Var}{\text{Var}}
\newcommand{\Lad}{\text{L}_{\text{ad}}^2}
\newcommand{\SI}[1]{\mathcal{I}\left[#1 \right]}
\newcommand{\FoxH}[5]{H_{#2}^{#1}\left(#3\:\middle\vert\: \begin{subarray}{l}#4\\[0.4em] #5\end{subarray}\right)}
% }}}

\newcommand*{\one}{{{\rm 1\mkern-1.5mu}\!{\rm I}}}

\begin{document}

\title{Answers to the referee report}

\author{Le Chen and Nicholas Eisenberg}
\date{\today}
\maketitle

The authors would like to thank the referee for the careful reading of the paper
and the valuable comments. We have revised the paper according to the referee's
report.

\bigskip\noindent\textbf{Answers to the six comments:}

\begin{enumerate}[(1)]
  \setlength{\itemsep}{10pt}

  \item[1.] page 3, paragraph following Definition 1.1:
    $\mathscr{L}(u(t,\cdot;\mu)$ is missing a parenthesis on the right
  \item[A:] We have added the missing right parenthesis for this appearance and
    a few others.


  \item[2.] page 5, the reviewer finds it instructive to view “Dalang’s
    function” $\Upsilon(\beta)$, for $\beta>0$, as being an interpolation
    between its strengthened counterparts, namely:
    \begin{align}\label{E:Interpolation}
      \Upsilon\left(\beta\right) \le 2 \Upsilon_\alpha(\beta)^{1/2} \Upsilon_\alpha(0)^{1/2}
    \end{align}
    which can be proven by simply splitting the integral into near and far
    fields, then optimizing the constant. This would complement the discussion
    on page 5/6, as well as the one in Section 3.2
  \item[A:] We thank the referee for bringing up this inequality with an effort
    to clarify the relation of the standard Dalang's condition and its
    strengthened versions. We agree with the referee and indeed,
    relation~\eqref{E:Interpolation} can be obtained by combining the following
    two inequalities:
    \begin{align*}
      \Upsilon\left(\beta\right) < \Upsilon\left(0\right)  \quad \text{and} \quad
      \Upsilon\left(\beta\right) < \beta^\alpha \Upsilon_\alpha\left(\beta\right).
    \end{align*}
    One can also obtain similar inequality such as
    \begin{align*}
      \Upsilon(\beta) < \min\left(\Upsilon\left(0\right),\; \beta^\alpha \Upsilon_\alpha\left(\beta\right)\right).
    \end{align*}
    In order to make the relation of the following four conditions clearer:
    \begin{align*}
      \Upsilon(\beta)<\infty, \quad
      \Upsilon(0)<\infty, \quad
      \Upsilon_\alpha(\beta)<\infty, \quad
      \Upsilon_\alpha(0)<\infty,
    \end{align*}
    we have made the following changes:
    \begin{enumerate}
      \item We have also included one footnote on p.~6 to provide one example constructed
        in~\cite{chen.huang.ea:19:dense} that
        \begin{align*}
          \Upsilon(\beta)<\infty, \quad
          \Upsilon(0)<\infty, \quad
          \Upsilon_\alpha(\beta)=\infty, \quad
          \Upsilon_\alpha(0)=\infty.
        \end{align*}
      \item We have slightly updated Remark~1.1.
      \item We have worked out Figure~1.1 with more information.
      \item We have expanded the discussion in Example~5.10 and included
        Figure~5.5 to show the subtlety of these conditions.
    \end{enumerate}

  \item[3.] page 7, proof of Corollary 1.2, second equation line: there is a ∗
    missing on $\zeta$ inside the last integral preceding the inequality. The
    integrand should be $(G(t,\cdot)*\rho)(y) |\zeta^*(y)|^2$
  \item[A:] It has been modified as suggested.

  \item[4.] page 9, definition of $H$: the authors should include the
    definition of $H$ since it is invoked several times in the paper and its
    properties rely on the restrictions associated to the Dalang-type condition.
    The reviewer went to reference [3] and does not think it will be difficult
    to add and will only add clarity to their framework.
  \item[A:] We have introduced Remark 2.1 to define $H$ and to state its basic
    properties.

  \item[5.] page 14, the convention in writing integrals is inconsistent.
    Although this point may seem pedantic, the authors write integrals in at
    least three different ways, for example as $\iint \ud x\ud y f (x, y) $ or
    $\iint f (x, y)\ud x\ud y$ or $\int\ud x (f x, y)\ud y$. The author
    identified that the main convention seems to be to write integrals as
    $\iint\ud x\ud y f(x,y)$ when calculations or estimates with integrals are
    being made, while $\iint f (x, y)\ud x\ud y$ when statements are asserted
    about them such as conditions being stated. The reviewer suggests that the
    authors make their convention uniform throughout the paper
  \item[A:] We have unified the notation for integral throughout the paper to
    the convention to be $\iint \ud x\ud y \: f(x,y)$. All integrals have been
    updated.

  \item[6.] page 26, sentence preceding Section 5.4: “Combing” should be
    “Combining”
  \item[A:] It has been modified as suggested.

\end{enumerate}

\bigskip\noindent\textbf{Additional changes that we have made}

\begin{enumerate}
  \item We have expressed our gratitude to the anonymous referee in the
    Acknowledgments.
  \item We have updated the funding information for the first author.
  \item The first case in Eq.~(5.21) has been updated. This is a copy-paste
    error during our edit.
  \item We have included the table of contents.
\end{enumerate}

\bibliographystyle{alpha}
\bibliography{All.bib}

\end{document}
